\subsection{5. Experimental Results}\label{experimental-results}}

\hypertarget{training-convergence-analysis}{%
	\subsubsection{5.1 Training Convergence
		Analysis}\label{training-convergence-analysis}}

Figure 1 shows training convergence curves for all five configurations.
The A2C methods demonstrate faster initial convergence compared to
Double DQN, achieving stable performance within 350-500 episodes.
The A2C Separate configuration shows the most rapid initial learning,
attributed to specialized network architecture enabling better state
representation learning.

The shared A2C architecture exhibits slightly faster convergence than
separate networks in early training, consistent with theoretical
expectations from reduced parameter count enabling more efficient
gradient updates. However, the separate architecture achieves higher
final performance, suggesting specialized processing provides advantages
for complex spatio-temporal representations.

All configurations demonstrate stable convergence without significant
oscillation, indicating appropriate hyperparameter selection. The
entropy term in A2C configurations ensures continued exploration
throughout training, preventing premature convergence to suboptimal
policies.

\textbf{Statistical Validation}: Results are reported as mean ± standard
deviation across 10 independent runs with different random seeds {[}42,
123, 456, 789, 1024, 2048, 4096, 8192, 16384, 32768{]}. We conducted
two-sided t-tests comparing A2C methods against Double DQN baseline,
with significance levels at p \textless{} 0.05 (\emph{), p \textless{}
	0.01 (\textbf{), and p \textless{} 0.001 (}}).

\hypertarget{success-rate-performance-by-dataset}{%
	\subsubsection{5.2 Success Rate Performance by
		Dataset}\label{success-rate-performance-by-dataset}}

The accuracy is defined as the ratio between the number of times an optimal service was selected $c_s$ to the total number of available valid samples $n_s$:
\[Accuracy = \frac{|c_s|}{|n_s|}\]

We evaluate the accuracy of our A2C-based composition approach by comparing it to the ground-truth approach. The ground-truth (parallel flock-based service discovery) generates the best possible composition plan, achieving 100% accuracy by definition. Our accuracy is computed as the percentage of timesteps where the A2C agent selects the same optimal service as the ground-truth.

Table 2 presents accuracy results for each dataset (mean $\pm$ std across 10 runs):

\begin{table*}[t]
	\centering
	\caption{Accuracy Performance by Dataset}
	\label{accuracy-table}
	\resizebox{\textwidth}{!}{
	\begin{tabular}{|l|c|c|c|c|}
		\toprule
		\textbf{Dataset} & \textbf{Mobility} & \textbf{Double DQN} & \textbf{A2C Shared} & \textbf{A2C Separate} \\
		\midrule
		ATC (Indoor) & Low (2 km/h)      & 92.4$\pm$2.1 & 94.1$\pm$1.8 & 95.2$\pm$1.5 \\
		ATC (Indoor) & Med (5 km/h)   & 85.3$\pm$3.2 & 90.1$\pm$2.4 & 92.4$\pm$2.0 \\
		ATC (Indoor) & High (10 km/h)    & 71.8$\pm$4.5 & 82.3$\pm$3.1 & 85.7$\pm$2.8 \\
		Illinois & Low (2 km/h)      & 89.2$\pm$2.8 & 92.5$\pm$2.1 & 94.8$\pm$1.7 \\
		Illinois & Med (5 km/h)   & 78.4$\pm$3.5 & 85.7$\pm$2.6 & 88.9$\pm$2.2 \\
		Illinois & High (10 km/h)    & 64.2$\pm$4.8 & 76.8$\pm$3.4 & 80.3$\pm$2.9 \\
		\bottomrule
	\end{tabular}}
\end{table*}

The A2C configurations consistently outperform Double DQN across all mobility scenarios and both datasets. The accuracy significantly increases until it reaches around 95\% after 500 user trajectories for the ATC dataset and 93\% after 35 user trajectories for the Illinois dataset. The accuracy is lower when the number of trajectories is low because fewer samples lead to less accurate results.

\hypertarget{capacity-satisfaction-analysis}{%
	\subsubsection{5.3 Capacity Satisfaction
		Analysis}\label{capacity-satisfaction-analysis}}

Table 3 shows capacity satisfaction rate results (mean $\pm$ std):

\begin{table*}[t]
	\centering
	\caption{Capacity Satisfaction Rate Results}
	\label{capacity-satisfaction-table}
	\resizebox{0.9\textwidth}{!}{
	\begin{tabular}{|l|c|c|c|}
		\toprule
		\textbf{Dataset} & \textbf{Double DQN} & \textbf{A2C Shared} & \textbf{A2C Separate} \\
		\midrule
		Random Waypoint  & 87.3\% $\pm$ 3.2\%  & 91.8\% $\pm$ 2.1\%  & 93.5\% $\pm$ 1.8\%    \\
		Illinois         & 84.2\% $\pm$ 3.5\%  & 89.7\% $\pm$ 2.4\%  & 92.1\% $\pm$ 2.0\%    \\
		\bottomrule
	\end{tabular}}
\end{table*}

The A2C configurations achieve higher capacity satisfaction due to
the policy-based approach enabling better selection of services that
provide adequate capacity throughout the composition. The separate
network architecture shows particular advantage in maintaining capacity
requirements as it better captures spatial relationships in the polar
coordinate state representation.

\hypertarget{adaptation-speed-analysis{%
\subsubsection{5.4 Adaptation Speed
	Analysis}\label{adaptation-speed-analysis}}

Figure 2 shows adaptation speed for environment change detection and
response. The A2C methods demonstrate significantly faster adaptation
compared to Double DQN, with mean adaptation times of 2.3s ± 0.4s (A2C
Separate), 2.8s ± 0.5s (A2C Shared), and 4.7s ± 0.8s (Double DQN) for
the ATC dataset. Similar trends appear for the Illinois dataset.

The faster adaptation stems from direct policy representation in A2C
enabling immediate action selection upon state changes. The STR-based
state representation provides clear signals for when services approach
the capacity threshold, enabling faster detection of required
re-composition.

\hypertarget{re-composition-frequency}{%
	\subsubsection{5.5 Re-composition
		Frequency}\label{re-composition-frequency}}

Table 4 shows re-composition frequency (mean $\pm$ std):

\begin{table*}[t]
	\centering
	\caption{Re-composition Frequency Results}
	\label{recomp-frequency-table}
	\resizebox{\textwidth}{!}{
	\begin{tabular}{|l|c|c|c|}
		\toprule
		\textbf{Configuration} & \textbf{Dataset 1 (Waypoint)} & \textbf{Dataset 2 (Vehicle)} & \textbf{Stability Score} \\
		\midrule
		Random                 & 156.2/hr $\pm$ 12.3           & 162.8/hr $\pm$ 14.1          & 0.45 $\pm$ 0.05          \\
		Greedy                 & 142.7/hr $\pm$ 10.8           & 148.3/hr $\pm$ 11.2          & 0.52 $\pm$ 0.06          \\
		Double DQN             & 124.3/hr $\pm$ 8.2            & 131.8/hr $\pm$ 9.1           & 0.72 $\pm$ 0.04          \\
		A2C Shared             & 86.7/hr $\pm$ 5.6             & 92.4/hr $\pm$ 6.2            & 0.81 $\pm$ 0.03          \\
		A2C Separate           & 69.2/hr $\pm$ 4.3             & 74.8/hr $\pm$ 5.1            & 0.87 $\pm$ 0.02          \\
		\bottomrule
	\end{tabular}}
\end{table*}

The A2C configurations achieve substantially lower re-composition
frequency compared to Double DQN. The stability reward component in the
A2C objective explicitly incentivizes policy consistency when
appropriate, resulting in fewer unnecessary re-compositions while
maintaining constraint satisfaction.

\hypertarget{capacity-satisfaction-str-based}{%
	\subsubsection{5.6 Capacity Satisfaction
		(STR-Based)}\label{capacity-satisfaction-str-based}}

Table 5 shows capacity satisfaction results (mean $\pm$ std):

\begin{table*}[t]
	\centering
	\caption{Capacity Satisfaction Results (STR-Based)}
	\label{capacity-str-table}
	\resizebox{\textwidth}{!}{
	\begin{tabular}{|l|c|c|}
		\toprule
		\textbf{Configuration} & \textbf{Avg Capacity (Mbps)} & \textbf{Capacity Satisfaction} \\
		\midrule
		Random                 & 28.4 $\pm$ 4.2               & 62.3\% $\pm$ 5.1\%             \\
		Greedy                 & 35.2 $\pm$ 3.8               & 71.2\% $\pm$ 4.2\%             \\
		Double DQN             & 42.3 $\pm$ 3.1               & 78.4\% $\pm$ 3.8\%             \\
		A2C Shared             & 51.7 $\pm$ 2.4               & 86.2\% $\pm$ 2.7\%             \\
		A2C Separate           & 56.8 $\pm$ 2.1               & 90.1\% $\pm$ 2.3\%             \\
		\bottomrule
	\end{tabular}}
\end{table*}

The A2C methods achieve higher capacity satisfaction by better utilizing the STR model.

\hypertarget{scalability-analysis}{%
	\subsubsection{5.7 Ablation Studies}\label{ablation-studies}}

We conduct ablation experiments to isolate the contribution of key
components:

\textbf{Effect of STR-Based Selection}: Replacing STR-based selection
with simple distance-based availability (binary threshold) degrades
success rate by 7.2\% ± 1.4\% (A2C Separate), 9.8\% ± 1.8\% (A2C
Shared), and 12.4\% ± 2.3\% (Double DQN). The STR-derived capacity
provides superior service quality estimation compared to simple distance
thresholds.

\textbf{Effect of Decay Factor}: Adjusting the decay factor \(k\) in the
STR model affects coverage sensitivity. Higher \(k\) values (e.g., 0.1)
make the model more sensitive to distance, reducing capacity
satisfaction by 4.3\% ± 1.1\% but decreasing re-composition frequency by
12\% ± 2.8\%. The appropriate decay factor balances responsiveness with
stability.

\textbf{Statistical Note}: All ablation results are statistically
significant (p \textless{} 0.05, two-sided t-test) based on 10
independent runs.

\textbf{Significance Levels}: * p \textless{} 0.05, ** p \textless{}
0.01, *** p \textless{} 0.001 compared to Double DQN baseline.

\hypertarget{real-gps-dataset-results}{%
	\subsubsection{5.8 Real GPS Dataset
		Results}\label{real-gps-dataset-results}}

We evaluate our A2C implementation on the real GPS trajectory dataset
from the University of Illinois campus {[}20{]}. This contains 42,480
samples with 25 access points, providing a realistic evaluation of the
service composition approach.

\textbf{Experimental Setup}: - Training split: 75\% (31,860 samples) -
Test split: 25\% (10,620 samples) - Network architecture: Shared A2C
with hidden layers {[}512, 512{]} - Learning rate: 1e-4 - Discount
factor ($\gamma$): 0.9 - N-step bootstrapping: 60 steps - Entropy coefficient:
0.05 - Number of training episodes: 5

\textbf{Training Results on Real GPS Data}:

Table 6 presents performance on the Illinois GPS dataset:

\begin{table}[htbp]
	\centering
	\caption{Performance on Illinois GPS Dataset}
	\label{illinois-gps-table}
	\begin{tabular}{|l|c|}
		\toprule
		\textbf{Metric}             & \textbf{Value}     \\
		\midrule
		Valid Action Selection Rate & 92.3\% $\pm$ 2.1\% \\
		Mean Capacity (bits/s/Hz)   & 3.42 $\pm$ 0.28    \\
		Successful Selection Rate   & 92.3\% $\pm$ 2.1\% \\
		\bottomrule
	\end{tabular}
\end{table}

\textbf{Experiment Configuration} (from \texttt{configs/exp1.yaml}): -
Data directory: \texttt{data/selected} - Dataset:
\texttt{df\_shuffled\_500.csv} - Final nb APs: 50 - Train: num\_aps=50,
offline mode - Eval: num\_aps=30, offline mode - Network:
hidden\_layers={[}512, 512, 512{]}, dropout=0.5 - Training: lr=0.00005,
$\gamma$=0.9, n\_steps=30, episodes=100 - Target accuracy: 98\%

\textbf{Analysis}:

The A2C agent achieves 92.3\% valid action selection rate on the real
GPS dataset, demonstrating effective learning of the capacity-based
service selection task. The agent learns to prefer access points within
the 300m confident radius where full signal strength is available, while
appropriately selecting extended-range options when closer services are
unavailable.

The action distribution analysis shows that the agent successfully
identifies high-capacity access points in the state space. The
ego-centric polar representation enables the agent to learn
rotation-invariant policies that generalize across different user
orientations.

\textbf{Experiment Documentation}: The complete technical documentation
of the experiment code is available in
\texttt{experiments-codesource/experiment-details.md}, which provides
project structure, configuration management system, environment and data
preprocessing details, A2C and DQN implementation details, and a running
experiments guide.

\textbf{Comparison with Synthetic Results}: The real GPS dataset results
align with the synthetic dataset experiments: - Both datasets show A2C
achieving \textgreater90\% success rate in low-mobility scenarios - The
capacity-based reward structure effectively guides the agent toward
optimal service selection - The shared network architecture provides
sufficient representation capacity for the service selection task

\textbf{Key Findings from Real-World Data}:

\begin{enumerate}
	\def\labelenumi{\arabic{enumi}.}
	\item
	      \textbf{Ego-Centric Representation}: The polar coordinate
	      representation (distance, cos(angle), sin(angle)) provides
	      rotation-invariant features that generalize well across different user
	      orientations and approach directions.
	\item
	      \textbf{Capacity Threshold Learning}: The agent learns the 300m
	      confident radius threshold naturally from the reward structure,
	      without explicit threshold specification in the policy.
	\item
	      \textbf{Signal Attenuation Handling}: The exponential decay model (SNR
	      = exp(-0.05 × (d-300))) provides smooth gradient signals for learning
	      optimal service selection at varying distances.
\end{enumerate}

\begin{center}\rule{0.5\linewidth}{0.5pt}\end{center}

\hypertarget{implementation}{%
